\documentclass[10pt,aspectratio=169,serif]{beamer}

\usepackage[fontset=none, scheme=plain]{ctex}
\IfFontExistsTF{Songti SC}{%
	\setCJKmainfont{Songti SC}[AutoFakeSlant=0.2]%
	\setCJKsansfont{Songti SC}[AutoFakeSlant=0.2]%
}{%
	\setCJKmainfont{FandolSong-Regular}[Extension=.otf,BoldFont=FandolSong-Bold,AutoFakeSlant=0.2]%
	\setCJKsansfont{FandolSong-Regular}[Extension=.otf,BoldFont=FandolSong-Bold]%
}
\usepackage{amsmath,amssymb,graphicx,tikz,booktabs,multicol,calligra}
\usepackage{pgfplots}
\pgfplotsset{compat=1.18}
\usepackage{tcolorbox}
\tcbuselibrary{skins}
\usepackage{XyygU}
\usepackage{caption}
\captionsetup[figure]{name=图,labelsep=period}

\usetikzlibrary{arrows.meta,positioning,calc}

% 公式间距
\setlength{\abovedisplayskip}{0pt plus 2pt minus 2pt}
\setlength{\belowdisplayskip}{0pt plus 2pt minus 2pt}

\title[诱导公式练习]{5.3 诱导公式（练习）}
\subtitle{公式一～公式六 \& “奇变偶不变，符号看象限”}
\author{}
\institute{}
\date{2026-10-16}

\begin{document}

\begin{frame}
  \titlepage
\end{frame}

% ========== 知识梳理：诱导公式总览 ==========
\begin{frame}[t]{诱导公式总览}
\small
\begin{tcolorbox}[colback=blue!5,colframe=blue,title=\textbf{六组诱导公式（$k\in\mathbf Z$）}]
\begin{enumerate}
  \item 公式一：$\sin(x+2k\pi)=\sin x$，$\cos(x+2k\pi)=\cos x$，$\tan(x+k\pi)=\tan x$；
  \item 公式二：$\sin(\pi+\alpha)=-\sin\alpha$，$\cos(\pi+\alpha)=-\cos\alpha$，$\tan(\pi+\alpha)=\tan\alpha$；
  \item 公式三：$\sin(-\alpha)=-\sin\alpha$，$\cos(-\alpha)=\cos\alpha$，$\tan(-\alpha)=-\tan\alpha$；
  \item 公式四：$\sin(\pi-\alpha)=\sin\alpha$，$\cos(\pi-\alpha)=-\cos\alpha$，$\tan(\pi-\alpha)=-\tan\alpha$；
  \item 公式五：$\sin\!\left(\dfrac{\pi}{2}-\alpha\right)=\cos\alpha$，$\cos\!\left(\dfrac{\pi}{2}-\alpha\right)=\sin\alpha$；
  \item 公式六：$\sin\!\left(\dfrac{\pi}{2}+\alpha\right)=\cos\alpha$，$\cos\!\left(\dfrac{\pi}{2}+\alpha\right)=-\sin\alpha$。
\end{enumerate}
\end{tcolorbox}
\vspace{0.3em}
\alert{口诀：}奇变偶不变，符号看象限。（把 $\alpha$ 看作锐角，看原函数名前角所在象限的符号）
\end{frame}

% ==================================================================
\section*{练习（一）公式一～公式四}
% ==================================================================

% ---------- 第1题 ----------
\begin{frame}[t]{练习一 第1题：化为锐角三角函数}
\textbf{1.} 将下列三角函数化为锐角三角函数，并填在题中横线上：
\begin{enumerate}[(1)]
  \item $\cos\dfrac{13\pi}{9}=$ \underline{\hspace{2cm}}；
  \item $\sin(1+\pi)=$ \underline{\hspace{2cm}}；
  \item $\sin\!\left(-\dfrac{\pi}{5}\right)=$ \underline{\hspace{2cm}}；
  \item $\tan(-70^\circ6')=$ \underline{\hspace{2cm}}；
  \item $\cos\dfrac{6\pi}{7}=$ \underline{\hspace{2cm}}；
  \item $\tan 1000^\circ21'=$ \underline{\hspace{2cm}}。
\end{enumerate}
\end{frame}

\begin{frame}[t]{练习一 第1题 解答}
\begin{tcolorbox}[colback=orange!10,colframe=orange,title=\textbf{解}]
先把角写成 $2k\pi+\theta$（或 $k\cdot360^\circ+\theta$）化到 $[0,2\pi)$（或 $[0,360^\circ)$），再用公式二～四化锐角。
\begin{enumerate}[(1)]
  \item $\dfrac{13\pi}{9}=\pi+\dfrac{4\pi}{9}$，故 $\cos\dfrac{13\pi}{9}=\alert{-\cos\dfrac{4\pi}{9}}$；
  \item $\sin(1+\pi)=\alert{-\sin 1}$；
  \item $\sin\!\left(-\dfrac{\pi}{5}\right)=\alert{-\sin\dfrac{\pi}{5}}$；
  \item $\tan(-70^\circ6')=\alert{-\tan 70^\circ6'}$；
  \item $\dfrac{6\pi}{7}=\pi-\dfrac{\pi}{7}$，故 $\cos\dfrac{6\pi}{7}=\alert{-\cos\dfrac{\pi}{7}}$；
  \item $1000^\circ21'=2\times360^\circ+280^\circ21'=2\times360^\circ+(360^\circ-79^\circ39')$，故 $\tan1000^\circ21'=\alert{-\tan79^\circ39'}$。
\end{enumerate}
\end{tcolorbox}
\end{frame}

% ---------- 第2题 ----------
\begin{frame}[t]{练习一 第2题：求三角函数值}
\textbf{2.} 利用公式求下列三角函数值：
\begin{enumerate}[(1)]
  \item $\cos(-420^\circ)$； \qquad
  \item $\sin\!\left(-\dfrac{7\pi}{6}\right)$； \qquad
  \item $\tan(-1140^\circ)$；
  \item $\cos\!\left(-\dfrac{77\pi}{6}\right)$； \qquad
  \item $\tan 315^\circ$； \qquad
  \item $\sin\!\left(-\dfrac{11\pi}{4}\right)$。
\end{enumerate}
\end{frame}

\begin{frame}[t]{练习一 第2题 解答}
\begin{tcolorbox}[colback=orange!10,colframe=orange,title=\textbf{解}]
\begin{enumerate}[(1)]
  \item $\cos(-420^\circ)=\cos420^\circ=\cos(360^\circ+60^\circ)=\cos60^\circ=\alert{\dfrac12}$；
  \item $\sin\!\left(-\dfrac{7\pi}{6}\right)=-\sin\dfrac{7\pi}{6}=-\sin\!\left(\pi+\dfrac{\pi}{6}\right)=\sin\dfrac{\pi}{6}=\alert{\dfrac12}$；
  \item $\tan(-1140^\circ)=-\tan1140^\circ=-\tan(3\times360^\circ+60^\circ)=-\tan60^\circ=\alert{-\sqrt3}$；
  \item $\cos\!\left(-\dfrac{77\pi}{6}\right)=\cos\dfrac{77\pi}{6}=\cos\!\left(12\pi+\dfrac{5\pi}{6}\right)=\cos\dfrac{5\pi}{6}=-\cos\dfrac{\pi}{6}=\alert{-\dfrac{\sqrt3}{2}}$；
  \item $\tan315^\circ=\tan(360^\circ-45^\circ)=-\tan45^\circ=\alert{-1}$；
  \item $\sin\!\left(-\dfrac{11\pi}{4}\right)=-\sin\dfrac{11\pi}{4}=-\sin\!\left(2\pi+\dfrac{3\pi}{4}\right)=-\sin\dfrac{3\pi}{4}=-\sin\dfrac{\pi}{4}=\alert{-\dfrac{\sqrt2}{2}}$。
\end{enumerate}
\end{tcolorbox}
\end{frame}

% ---------- 第3题 ----------
\begin{frame}[t]{练习一 第3题：化简}
\textbf{3.} 化简：
\begin{enumerate}[(1)]
  \item $\sin(-\alpha-180^\circ)\cos(-\alpha)\sin(-\alpha+180^\circ)$；
  \item $\cos^3(-\alpha)\sin(2\pi+\alpha)\tan^3(-\alpha-\pi)$。
\end{enumerate}
\end{frame}

\begin{frame}[t]{练习一 第3题 解答}
\begin{tcolorbox}[colback=orange!10,colframe=orange,title=\textbf{解}]
(1) 逐项化简：
\[
\sin(-\alpha-180^\circ)=-\sin(\alpha+180^\circ)=-(-\sin\alpha)=\sin\alpha,
\]
\[
\cos(-\alpha)=\cos\alpha,\qquad \sin(-\alpha+180^\circ)=\sin(180^\circ-\alpha)=\sin\alpha.
\]
故原式 $=\sin\alpha\cdot\cos\alpha\cdot\sin\alpha=\alert{\sin^2\alpha\cos\alpha}$。
\vspace{0.4em}
(2) $\cos(-\alpha)=\cos\alpha$，$\sin(2\pi+\alpha)=\sin\alpha$，
\[
\tan(-\alpha-\pi)=-\tan(\alpha+\pi)=-\tan\alpha.
\]
故原式 $=\cos^3\alpha\cdot\sin\alpha\cdot(-\tan^3\alpha)
=-\cos^3\alpha\sin\alpha\cdot\dfrac{\sin^3\alpha}{\cos^3\alpha}
=\alert{-\sin^4\alpha}$。
\end{tcolorbox}
\end{frame}

% ---------- 第4题 ----------
\begin{frame}[t]{练习一 第4题：填表}
\textbf{4.} 填表：
\vspace{0.5em}
\begin{center}
\begin{tabular}{c|cccccc}
$\alpha$ & $-\dfrac{4\pi}{3}$ & $-\dfrac{5\pi}{4}$ & $-\dfrac{5\pi}{3}$ & $-\dfrac{7\pi}{4}$ & $-\dfrac{8\pi}{3}$ & $-\dfrac{11\pi}{4}$\\
\hline
$\sin\alpha$ & & & & & & \\
\hline
$\cos\alpha$ & & & & & & \\
\hline
$\tan\alpha$ & & & & & & \\
\end{tabular}
\end{center}
\end{frame}

\begin{frame}[t]{练习一 第4题 解答}
\begin{tcolorbox}[colback=orange!10,colframe=orange,title=\textbf{解}]
把各角加上 $2\pi$ 或 $4\pi$ 化到 $[0,2\pi)$：
\[
-\tfrac{4\pi}{3}\to\tfrac{2\pi}{3},\ -\tfrac{5\pi}{4}\to\tfrac{3\pi}{4},\ -\tfrac{5\pi}{3}\to\tfrac{\pi}{3},\ -\tfrac{7\pi}{4}\to\tfrac{\pi}{4},\ -\tfrac{8\pi}{3}\to\tfrac{4\pi}{3},\ -\tfrac{11\pi}{4}\to\tfrac{5\pi}{4}.
\]
\vspace{0.3em}
\small
\begin{tabular}{c|cccccc}
$\alpha$ & $-\tfrac{4\pi}{3}$ & $-\tfrac{5\pi}{4}$ & $-\tfrac{5\pi}{3}$ & $-\tfrac{7\pi}{4}$ & $-\tfrac{8\pi}{3}$ & $-\tfrac{11\pi}{4}$\\
\hline
$\sin\alpha$ & $\tfrac{\sqrt3}{2}$ & $\tfrac{\sqrt2}{2}$ & $\tfrac{\sqrt3}{2}$ & $\tfrac{\sqrt2}{2}$ & $-\tfrac{\sqrt3}{2}$ & $-\tfrac{\sqrt2}{2}$\\
\hline
$\cos\alpha$ & $-\tfrac12$ & $-\tfrac{\sqrt2}{2}$ & $\tfrac12$ & $\tfrac{\sqrt2}{2}$ & $-\tfrac12$ & $-\tfrac{\sqrt2}{2}$\\
\hline
$\tan\alpha$ & $-\sqrt3$ & $-1$ & $\sqrt3$ & $1$ & $\sqrt3$ & $1$\\
\end{tabular}
\end{tcolorbox}
\end{frame}

% ==================================================================
\section*{练习（二）公式五～公式六}
% ==================================================================

% ---------- 第1题 ----------
\begin{frame}[t]{练习二 第1题：求三角函数值}
\textbf{1.} 用诱导公式求下列三角函数值（可用计算工具，第(3)(4)(6)题精确到 $0.0001$）：
\begin{enumerate}[(1)]
  \item $\cos\dfrac{65\pi}{6}$； \qquad
  \item $\sin\!\left(-\dfrac{31\pi}{4}\right)$； \qquad
  \item $\cos(-1182^\circ13')$；
  \item $\sin 670^\circ39'$； \qquad
  \item $\tan\!\left(-\dfrac{26\pi}{3}\right)$； \qquad
  \item $\tan 580^\circ21'$。
\end{enumerate}
\end{frame}

\begin{frame}[t]{练习二 第1题 解答}
\begin{tcolorbox}[colback=orange!10,colframe=orange,title=\textbf{解}]
\begin{enumerate}[(1)]
  \item $\dfrac{65\pi}{6}=10\pi+\dfrac{5\pi}{6}$，$\cos\dfrac{65\pi}{6}=\cos\dfrac{5\pi}{6}=-\dfrac{\sqrt3}{2}\approx\alert{-0.8660}$；
  \item $\sin\!\left(-\dfrac{31\pi}{4}\right)=-\sin\!\left(6\pi+\dfrac{7\pi}{4}\right)=-\sin\dfrac{7\pi}{4}=\dfrac{\sqrt2}{2}\approx\alert{0.7071}$；
  \item $\cos(-1182^\circ13')=\cos1182^\circ13'=\cos(3\times360^\circ+102^\circ13')=\cos102^\circ13'=-\cos77^\circ47'\approx\alert{-0.2116}$；
  \item $\sin670^\circ39'=\sin(360^\circ+310^\circ39')=\sin(360^\circ-49^\circ21')=-\sin49^\circ21'\approx\alert{-0.7587}$；
  \item $\tan\!\left(-\dfrac{26\pi}{3}\right)=-\tan\!\left(8\pi+\dfrac{2\pi}{3}\right)=-\tan\dfrac{2\pi}{3}=\sqrt3\approx\alert{1.7321}$；
  \item $\tan580^\circ21'=\tan(360^\circ+220^\circ21')=\tan(180^\circ+40^\circ21')=\tan40^\circ21'\approx\alert{0.8496}$。
\end{enumerate}
\end{tcolorbox}
\end{frame}

% ---------- 第2题 ----------
\begin{frame}[t]{练习二 第2题：证明}
\textbf{2.} 证明：
\begin{enumerate}[(1)]
  \item $\cos\!\left(\dfrac{5\pi}{2}-\alpha\right)=\sin\alpha$； \qquad
  \item $\cos\!\left(\dfrac{7\pi}{2}+\alpha\right)=\sin\alpha$；
  \item $\sin\!\left(\dfrac{9\pi}{2}-\alpha\right)=\cos\alpha$； \qquad
  \item $\sin\!\left(\dfrac{11\pi}{2}-\alpha\right)=-\cos\alpha$。
\end{enumerate}
\end{frame}

\begin{frame}[t]{练习二 第2题 解答}
\begin{tcolorbox}[colback=orange!10,colframe=orange,title=\textbf{证明}]
(1) $\cos\!\left(\dfrac{5\pi}{2}-\alpha\right)=\cos\!\left(2\pi+\dfrac{\pi}{2}-\alpha\right)=\cos\!\left(\dfrac{\pi}{2}-\alpha\right)=\alert{\sin\alpha}$。
\vspace{0.3em}
(2) $\cos\!\left(\dfrac{7\pi}{2}+\alpha\right)=\cos\!\left(3\pi+\dfrac{\pi}{2}+\alpha\right)=\cos\!\left(\pi+\dfrac{\pi}{2}+\alpha\right)=-\cos\!\left(\dfrac{\pi}{2}+\alpha\right)=-(-\sin\alpha)=\alert{\sin\alpha}$。
\vspace{0.3em}
(3) $\sin\!\left(\dfrac{9\pi}{2}-\alpha\right)=\sin\!\left(4\pi+\dfrac{\pi}{2}-\alpha\right)=\sin\!\left(\dfrac{\pi}{2}-\alpha\right)=\alert{\cos\alpha}$。
\vspace{0.3em}
(4) $\sin\!\left(\dfrac{11\pi}{2}-\alpha\right)=\sin\!\left(4\pi+\dfrac{3\pi}{2}-\alpha\right)=\sin\!\left(\pi+\left(\dfrac{\pi}{2}-\alpha\right)\right)=-\sin\!\left(\dfrac{\pi}{2}-\alpha\right)=\alert{-\cos\alpha}$。
\end{tcolorbox}
\end{frame}

% ---------- 第3题 ----------
\begin{frame}[t]{练习二 第3题：化简}
\textbf{3.} 化简：
\begin{enumerate}[(1)]
  \item $\dfrac{\cos\!\left(\alpha-\frac{\pi}{2}\right)}{\sin\!\left(\frac{5\pi}{2}+\alpha\right)}\sin(\alpha-2\pi)\cos(2\pi-\alpha)$；
  \item $\cos^2(-\alpha)-\dfrac{\tan(2\pi+\alpha)}{\cos\!\left(\frac{\pi}{2}+\alpha\right)}$；
  \item $\dfrac{\cos(\alpha-3\pi)\cos\!\left(\frac{3\pi}{2}-\alpha\right)}{\sin^2\!\left(\alpha-\frac{\pi}{2}\right)}$。
\end{enumerate}
\end{frame}

\begin{frame}[t]{练习二 第3题 解答}
\begin{tcolorbox}[colback=orange!10,colframe=orange,title=\textbf{解}]
(1) $\cos\!\left(\alpha-\dfrac{\pi}{2}\right)=\cos\!\left(\dfrac{\pi}{2}-\alpha\right)=\sin\alpha$，
$\sin\!\left(\dfrac{5\pi}{2}+\alpha\right)=\sin\!\left(\dfrac{\pi}{2}+\alpha\right)=\cos\alpha$，
$\sin(\alpha-2\pi)=\sin\alpha$，$\cos(2\pi-\alpha)=\cos\alpha$。
故原式 $=\dfrac{\sin\alpha}{\cos\alpha}\cdot\sin\alpha\cdot\cos\alpha=\alert{\sin^2\alpha}$。
\vspace{0.3em}
(2) $\cos^2(-\alpha)=\cos^2\alpha$，$\tan(2\pi+\alpha)=\tan\alpha$，$\cos\!\left(\dfrac{\pi}{2}+\alpha\right)=-\sin\alpha$。
故原式 $=\cos^2\alpha-\dfrac{\tan\alpha}{-\sin\alpha}=\cos^2\alpha+\dfrac{1}{\cos\alpha}=\alert{\cos^2\alpha+\sec\alpha}$。
\vspace{0.3em}
(3) $\cos(\alpha-3\pi)=\cos(3\pi-\alpha)=\cos(\pi-\alpha)=-\cos\alpha$，
$\cos\!\left(\dfrac{3\pi}{2}-\alpha\right)=-\sin\alpha$，
$\sin^2\!\left(\alpha-\dfrac{\pi}{2}\right)=\sin^2\!\left(\dfrac{\pi}{2}-\alpha\right)=\cos^2\alpha$。
故原式 $=\dfrac{(-\cos\alpha)(-\sin\alpha)}{\cos^2\alpha}=\alert{\tan\alpha}$。
\end{tcolorbox}
\end{frame}

\begin{frame}[c]
  \begin{center}
    {\Huge\calligra Thanks!}
  \end{center}
\end{frame}

\end{document}
